% BOOK-FILE-ID: f-847298b7c959
\subsection{変数係数の一次型}\label{節：変数係数の一次型}
階差等比複合型$a_{n+1}=pa_n+f(n)$では,掛ける数$p$は定数であった.
では,この$p$までもが$n$によって変わったらどうなるだろうか.
一般的な形は
\[
  a_{n+1}=p_na_n+q_n
\]
で与えられる.\,$p_n,q_n$は$n$だけで決まる数である.
階比型と階差型が組み合わさったような漸化式なので,その2つを使って既知の形に帰着させることを目標にしよう.

\subsubsection{例題：両辺を割って階差型にする}
まずは具体的な式で考えよう.
\BookMathLayout{\[
    a_1=2,\qquad
    a_{n+1}=\frac{n+2}{n}a_n+(n+1)(n+2)
    \quad(n\ge1)
  \]}{\[
    \begin{aligned}
      a_1&=2,\\
      a_{n+1}&=\frac{n+2}{n}a_n+(n+1)(n+2)\\
      &\quad(n\ge1)
    \end{aligned}
  \]}
を考える.
邪魔なのは,\,$a_n$に掛かっている$\frac{n+2}{n}$である.
これさえなければ階差型として解ける.
\sectionref*{節：階比型}では,両辺で同じ形を作るために,うまい式で割ったことを思い出そう.
係数の分母$n$と分子$n+2$の間には$n+1$が抜けているので,
\[
  \frac{(n+1)(n+2)}{n(n+1)}=\frac{n+2}{n}
\]
と書ける.つまり,\,$n(n+1)$の添字を1つ進めると,ちょうど$\frac{n+2}{n}$倍になるのだ.
そこで両辺を$(n+1)(n+2)$で割ってやると,
\[
  \frac{a_{n+1}}{(n+1)(n+2)}
  =\frac{a_n}{n(n+1)}+1
\]
となり,左辺が$n+1$に関する式,右辺が$n$に関する式という"きれい"な形になった.
ここで$b_n=\frac{a_n}{n(n+1)}$と置くと,
\[
  b_{n+1}=b_n+1,\qquad b_1=\frac{a_1}{1\cdot2}=1
\]
という等差型になるので,\,$b_n=n$である.
したがって
\[
  a_n=n^2(n+1)
\]
を得る.
初項を一般の$a_1=s$に変えても,\,$b_1=\frac{s}{2}$から同じように
\[
  a_n=n(n+1)\left(n+\frac{s-2}{2}\right)
\]
となる.

\subsubsection{一般の式への拡張}
同じことを一般の$a_{n+1}=p_na_n+q_n$で行ってみよう.
先ほどの$n(n+1)$の役割は,「添字を1つ進めると$p_n$倍になる数列」であった.
そこで
\[
  h_1=1,\qquad h_{n+1}=p_nh_n
\]
となる数列$h_n$を用意する.これは階比型そのものなので,
\[
  h_n=\prod_{k=1}^{n-1}p_k
\]
である.\footnote{何も掛けられていない$n=1$のときは$1$とする.また,ここではすべての$p_n$が0でないとする.もし$p_j=0$なら,\,$a_{j+1}=q_j$はそれ以前の項によらず決まるので,これを新たな初項として考え直せばよい.}
元の式を$h_{n+1}$で割ると,\footnote{このように,階比型の解$h_n$で割って$a_n$の係数を1にする操作を正規化という.詳しくは付録の\sectionref{節：正規化}で扱う.}
\[
  \frac{a_{n+1}}{h_{n+1}}=\frac{a_n}{h_n}+\frac{q_n}{h_{n+1}}
\]
となり,\,$b_n=\frac{a_n}{h_n}$は階差型になる.\,$b_1=a_1$から書き下していけば
\[
  b_n=a_1+\sum_{k=1}^{n-1}\frac{q_k}{h_{k+1}}
\]
であり,\,$a_n=h_nb_n$に戻して
\[
  a_n=h_n\left(a_1+\sum_{k=1}^{n-1}\frac{q_k}{h_{k+1}}\right)
\]
を得る.\footnote{先ほどの例題では$h_n=\frac{n(n+1)}{2}$だが,定数倍した$n(n+1)$で割っても,初項$b_1$を正しく合わせれば同じ答えになる.}
ただし,これは一般項を積と和で表しただけである.
その積や和が最後まで計算できるかどうかは,\,$p_n,q_n$次第なのだ.

\begin{bookpoint}
  \textbf{POINT}\par
  $h_{n+1}=p_nh_n$となる$h_n$で両辺を割り,\,$\frac{a_n}{h_n}$についての階差型に帰着させる.
\end{bookpoint}

\subsubsection{別解：同じ漸化式を満たす式を引く}
ここまでは係数$p_n$を消すことを考えてきたが,付加項$q_n$の方を消すこともできる.
\sectionref*{節：階差等比複合型}で不動関数を引いたのと同じ発想である.
先ほどの例題で得た$v_n=n^2(n+1)$は,元の漸化式を満たす.
そこで$d_n=a_n-v_n$と置き,
\[
  \begin{array}{rrl}
    &a_{n+1}={}&\frac{n+2}{n}a_n+(n+1)(n+2)\\
    -\,)&v_{n+1}={}&\frac{n+2}{n}v_n+(n+1)(n+2)\\ \hline
    &d_{n+1}={}&\frac{n+2}{n}d_n
  \end{array}
\]
とすれば付加項が消え,階比型になる.初項$a_1=s$なら$d_1=s-2$であり,
\[
  d_n=\frac{s-2}{2}n(n+1),\qquad a_n=v_n+d_n
\]
から同じ答えを得る.
一般項が「元の漸化式を満たす1つの式」と「付加項のない漸化式を満たす式の定数倍」の和になっていることに注意しよう.
特性方程式型でも階差等比複合型でも,同じ形の一般項が現れていたはずである.
なぜいつもこの形になるのかは,\sectionref*{節：答えの形を見通す}で考える.

\begin{bookpoint}
  \textbf{POINT}\par
  係数$p_n$を消したければ$h_{n+1}=p_nh_n$となる$h_n$で割り,付加項$q_n$を消したければ同じ漸化式を満たす式を引く.
\end{bookpoint}
操作だけを比べる練習は\sectionref*{節：操作を選ぶ練習}で,
\bookterm{general}{一般項}まで求める練習は\bookproblemref{practice-connections}{1}{接続演習問1}で行う.
